\documentclass[12pt,a4paper]{extarticle}
\usepackage[utf8]{inputenc}
\usepackage[T5]{fontenc}
\usepackage{geometry}
\geometry{a4paper, left=1.5cm, right=1cm, top=1.5cm, bottom=1.5cm, headheight=20pt, headsep=15pt, footskip=28pt}
\usepackage{enumitem}
% Gọi package nằm trong thư mục con template
\usepackage{../template/mngocbox}

\begin{document}

% Sử dụng ngoặc vuông [...] để thay đổi linh hoạt các thông số
\makeheadbox[headerright={Tổng ôn 10-11},chude={CHỦ ĐỀ 01},tieude={Vi-et},dinhdang={Đại số},lop={12 },gv={Nguyễn Lê Minh Ngọc}]

\vspace{2em}
\makeheading{IV}{Đề Kiểm tra}



\begin{enumerate}
    \item[\textbf{Câu 1:}] Cho phương trình $x^2 + mx + 2m - 4 = 0$. Tìm $m$ để phương trình có hai nghiệm phân biệt $x_1, x_2$ thỏa mãn $|x_1| + |x_2| = 3$.
    
    \item[\textbf{Câu 2:}] Cho phương trình $x^2 - 4x - m^2 - 1 = 0$. Tìm $m$ để phương trình có hai nghiệm phân biệt $x_1, x_2$ phân biệt thỏa mãn $x_2 = -5x_1$.
    
    \item[\textbf{Câu 3:}] Cho phương trình $x^2 - 2(k-1)x - 4k = 0$. Tìm $k$ để phương trình có hai nghiệm phân biệt $x_1, x_2$ phân biệt thỏa mãn $3x_1 - x_2 = 2$.
    
    \item[\textbf{Câu 4:}] Cho phương trình $x^2 - 6x + m + 3 = 0$. Tìm $m$ để phương trình có hai nghiệm phân biệt $x_1, x_2$ phân biệt thỏa mãn $x_2 = x_1^2$.
    
    \item[\textbf{Câu 5:}] Cho phương trình $x^2 - (m-3)x - 5 = 0$. Tìm $m$ để phương trình có hai nghiệm $x_1, x_2$ là các số nguyên.
    
    \item[\textbf{Câu 6:}] Cho phương trình $2x^2 - 3x - 1 = 0$ có hai nghiệm là $x_1, x_2$, không giải phương trình hãy tính giá trị của biểu thức $A = \frac{x_1 - 1}{x_2 + 1} + \frac{x_2 - 1}{x_1 + 1}$.
    
    \item[\textbf{Câu 7:}] Cho phương trình: $x^2 - 5x - 6 = 0$ có hai nghiệm $x_1, x_2$. Hãy tính giá trị của biểu thức sau: $A = \frac{x_1}{x_2 - 1} + \frac{x_2}{x_1 - 1}$.
    
    \item[\textbf{Câu 8:}] Cho phương trình: $x^2 + 3x - 10 = 0$ có 2 nghiệm $x_1, x_2$. Tính giá trị biểu thức $A = \frac{x_1 + 2}{x_2} + \frac{x_2 + 2}{x_1}$.
    
    \item[\textbf{Câu 9:}] Gọi $x_1, x_2$ là hai nghiệm của phương trình: $x^2 - x - 1 = 0$. Lập phương trình bậc hai có hai nghiệm là $\frac{x_2 + 1}{x_1}, \frac{x_1 + 1}{x_2}$.
    
    \item[\textbf{Câu 10:}] Cho phương trình $3x^2 - 6x + 2 = 0$ có hai nghiệm là $x_1, x_2$. Không giải phương trình, hãy tính giá trị của biểu thức $A = x_1^2 + x_2^2 - x_1x_2$.
\end{enumerate}

\end{document}